\chapter{Аналитическая часть}
%<пишем про тему и все, связанное с ней>

\section{Потоки}

Поток -- наименьшая единица работы, которая может быть назначена операционной системе для выполнения. Поток выполняется внутри процесса, и несколько потоков могут существовать внутри одного процесса, работая параллельно~\cite{lit3}. 

Нативный поток -- поток выполнения, которым управляет непосредственно операционная система и которые напрямую взаимодействуют с ядром операционной системы. 

Зеленый поток -- поток, который управляется средой выполнения (например, виртуальной машиной).

\section{Связность графа}

Связность графа -- качество графа, отображающее наличие путей между его вершинами и их количество~\cite{lit2}.

Компонента связности -- максимальный подграф неориентированного графа, вершины которого соединены между собой. 
Компонента сильной связности -- максимальный подграф ориентированного графа, каждая пара вершин которого достижима между собой.
Компонента слабой связности -- максимальный подграф ориентированного графа, каждая пара вершин которого достижима между собой, если этот граф рассматривается как неориентированный.

\section{Задание по варианту}

Проверка связности графа, поиск всех компонент связности (ответ – по списку вершин на каждую компоненту связности).

Дан граф в формате graphviz. Входные данные считываются из файла. Граф не полносвязный, количество вершин не менее 3000, из/в вершину ведут, вероятнее всего, от 0 до 15 дуг.

\section*{Вывод}
%<что сделали в аналите, кратко>

В данном разделе мной был проведен анализ теоретических материалов, объясняющих принцип работы алгоритмов.

\clearpage
